% GENERATED FILE. Edit canonical lessons, then run npm run generate:books.
% canonical-source-hash: fnv1a64:b4d21c90ae24fbef
% canonical-lessons: FR-C407-stationnement, FR-C407-stationner, FR-C407-horodateur, FR-C407-creneau, FR-C407-fourriere, FR-R407-revision-400-407

\chapter{Se garer en ville}
\label{ch:fr-se-garer-en-ville}
\hlchaptermodality{407}

\begin{chapteropening}
\textbf{By the end of this chapter:} \emph{I can talk about parking in town: parking, stopping, the ticket machine, a tight parking space and the pound.}

Complete the last lesson of chapter 407: I can talk about parking in town: parking, stopping, the ticket machine, a tight parking space and the pound.
\end{chapteropening}

\section[le stationnement]{le stationnement — parking}

\label{lesson:FR-C407-stationnement}



\begin{quote}
Before the new one: say the French for the boot, then the French for the back seat.
\end{quote}

\subsection*{You'll want to know: le stationnement}
\textbf{le stationnement} — \textquotedblleft{}parking\textquotedblright{}.

\begin{quote}
Le stationnement est interdit devant l'école.
\end{quote}

\begin{quote}
\emph{Parking is not allowed in front of the school.}
\end{quote}

\emph{Stationnement payant} — pay parking.

\begin{cousinweb}
\emph{stationnement} is built on \emph{stationner} 'to park' with \emph{-ment}. \emph{Stationner} comes from \emph{station}, a learned borrowing from Latin \emph{statiō} 'a standing', from \emph{stāre} 'to stand' (EtymDB).
\end{cousinweb}

\begin{grammarlens}[title={What you've built}]
One more word for telling, explaining and describing at B1.
\end{grammarlens}

\subsection*{Your turn}
\begin{itemize}
\raggedright
  \item \emph{Say it:} \emph{le stationnement}
  \item \emph{Say it:} \emph{le stationnement}, once more
  \item \emph{Say it:} a sentence of your own with \emph{le stationnement}
  \item \emph{Recall:} say the French for the boot, then the French for the back seat, then say \emph{le stationnement} again
\end{itemize}

\begin{tcolorbox}[breakable,colback=teal!4,colframe=teal!35!black,title={Before you move on}]
What does le stationnement mean? (\textquotedblleft{}Parking\textquotedblright{}.) Say it once more.
\end{tcolorbox}
\section[stationner]{stationner — to park, to be parked}

\label{lesson:FR-C407-stationner}



\begin{quote}
Before the new one: say the French for the back seat, then the French for parking.
\end{quote}

\subsection*{You'll want to know: stationner}
\textbf{stationner} — \textquotedblleft{}to park, to be parked\textquotedblright{}.

\begin{quote}
Les camions ne peuvent pas stationner ici.
\end{quote}

\begin{quote}
\emph{Lorries cannot park here.}
\end{quote}

More formal than \emph{se garer}.

\begin{cousinweb}
\emph{stationner} is built on \emph{station}, a learned borrowing from Latin \emph{statiō} 'a standing', from \emph{stāre} 'to stand' (EtymDB).
\end{cousinweb}

\begin{grammarlens}[title={What you've built}]
One more word for telling, explaining and describing at B1.
\end{grammarlens}

\subsection*{Your turn}
\begin{itemize}
\raggedright
  \item \emph{Say it:} \emph{stationner}
  \item \emph{Say it:} \emph{stationner}, once more
  \item \emph{Say it:} a sentence of your own with \emph{stationner}
  \item \emph{Recall:} say the French for the back seat, then the French for parking, then say \emph{stationner} again
\end{itemize}

\begin{tcolorbox}[breakable,colback=teal!4,colframe=teal!35!black,title={Before you move on}]
What does stationner mean? (\textquotedblleft{}To park, to be parked\textquotedblright{}.) Say it once more.
\end{tcolorbox}
\section[l'horodateur]{l'horodateur — the parking ticket machine}

\label{lesson:FR-C407-horodateur}



\begin{quote}
Before the new one: say the French for parking, then the French for to park.
\end{quote}

\subsection*{You'll want to know: l'horodateur}
\textbf{l'horodateur} — \textquotedblleft{}the parking ticket machine\textquotedblright{}.

\begin{quote}
Elle paie à l'horodateur avec sa carte.
\end{quote}

\begin{quote}
\emph{She pays at the ticket machine with her card.}
\end{quote}

A masculine noun. Older single meters were \emph{parcmètres}.

\begin{cousinweb}
\emph{horodateur} is built from \emph{horo-}, from Greek \emph{h\'{\={o}}ra} 'hour', and \emph{dateur} 'something that dates', from \emph{date}: it prints the time on the ticket. EtymDB does not list it.
\end{cousinweb}

\begin{grammarlens}[title={What you've built}]
One more word for telling, explaining and describing at B1.
\end{grammarlens}

\subsection*{Your turn}
\begin{itemize}
\raggedright
  \item \emph{Say it:} \emph{l'horodateur}
  \item \emph{Say it:} \emph{l'horodateur}, once more
  \item \emph{Say it:} a sentence of your own with \emph{l'horodateur}
  \item \emph{Recall:} say the French for parking, then the French for to park, then say \emph{l'horodateur} again
\end{itemize}

\begin{tcolorbox}[breakable,colback=teal!4,colframe=teal!35!black,title={Before you move on}]
What does l'horodateur mean? (\textquotedblleft{}The parking ticket machine\textquotedblright{}.) Say it once more.
\end{tcolorbox}
\section[le créneau]{le créneau — the parallel parking space; the time slot}

\label{lesson:FR-C407-creneau}



\begin{quote}
Before the new one: say the French for to park, then the French for the parking ticket machine.
\end{quote}

\subsection*{You'll want to know: le créneau}
\textbf{le créneau} — \textquotedblleft{}the parallel parking space; the time slot\textquotedblright{}.

\begin{quote}
Mon frère fait très bien les créneaux.
\end{quote}

\begin{quote}
\emph{My brother is very good at parallel parking.}
\end{quote}

\emph{Faire un créneau} — to parallel park. Also a time slot: \emph{un créneau libre}.

\begin{cousinweb}
\emph{créneau} comes from Old French \emph{crenel}, a small \emph{cren} 'notch': first the gap at the top of a castle wall. The origin of \emph{cren} is not settled.
\end{cousinweb}

\begin{grammarlens}[title={What you've built}]
One more word for telling, explaining and describing at B1.
\end{grammarlens}

\subsection*{Your turn}
\begin{itemize}
\raggedright
  \item \emph{Say it:} \emph{le créneau}
  \item \emph{Say it:} \emph{le créneau}, once more
  \item \emph{Say it:} a sentence of your own with \emph{le créneau}
  \item \emph{Recall:} say the French for to park, then the French for the parking ticket machine, then say \emph{le créneau} again
\end{itemize}

\begin{tcolorbox}[breakable,colback=teal!4,colframe=teal!35!black,title={Before you move on}]
What does le créneau mean? (\textquotedblleft{}The parallel parking space; the time slot\textquotedblright{}.) Say it once more.
\end{tcolorbox}
\section[la fourrière]{la fourrière — the car pound}

\label{lesson:FR-C407-fourriere}



\begin{quote}
Before the new one: say the French for the parking ticket machine, then the French for the parallel parking space.
\end{quote}

\subsection*{You'll want to know: la fourrière}
\textbf{la fourrière} — \textquotedblleft{}the car pound\textquotedblright{}.

\begin{quote}
Sa voiture était mal garée et elle est partie à la fourrière.
\end{quote}

\begin{quote}
\emph{Her car was badly parked and it was taken to the pound.}
\end{quote}

Also the pound for stray dogs.

\begin{cousinweb}
\emph{fourrière} is usually traced to Old French \emph{feurre} 'fodder', of Germanic origin: first the place where seized animals were kept and fed. EtymDB does not list it.
\end{cousinweb}

\begin{grammarlens}[title={What you've built}]
One more word for telling, explaining and describing at B1.
\end{grammarlens}

\subsection*{Your turn}
\begin{itemize}
\raggedright
  \item \emph{Say it:} \emph{la fourrière}
  \item \emph{Say it:} \emph{la fourrière}, once more
  \item \emph{Say it:} a sentence of your own with \emph{la fourrière}
  \item \emph{Recall:} say the French for the parking ticket machine, then the French for the parallel parking space, then say \emph{la fourrière} again
\end{itemize}

\begin{tcolorbox}[breakable,colback=teal!4,colframe=teal!35!black,title={Before you move on}]
What does la fourrière mean? (\textquotedblleft{}The car pound\textquotedblright{}.) Say it once more.
\end{tcolorbox}
\section[(dialogue)]{Routes, gares et hôtels}

\label{lesson:FR-R407-revision-400-407}



\begin{quote}
Say the French for the ring road and the outskirts.
\end{quote}

\subsection*{Your turn}
Say these aloud:

\begin{itemize}
\raggedright
  \item the ring road, the outskirts, the interchange, the slip road, to turn off
  \item to speed up, to turn the wheel, to skid, to stall, the pedal
  \item the windscreen, the rear-view mirror, the horn, to honk, to mist up
  \item the pump, the diesel, the fuel, the tank, to start up
  \item the puncture, to get a flat tyre, the tyre, the jack, to pump up
  \item to get going again, the tow truck, the breakdown mechanic, to tow, the garage owner
  \item the bodywork, the bonnet, the car door, the boot, the back seat
  \item parking, to park, the parking ticket machine, the parallel parking space, the car pound
\end{itemize}

\begin{tcolorbox}[breakable,colback=teal!4,colframe=teal!35!black,title={Before you move on}]
The ring road? (\textbf{le périphérique}.) The outskirts? (\textbf{la périphérie}.) The interchange? (\textbf{l'échangeur}.) The slip road? (\textbf{la bretelle}.) To turn off? (\textbf{bifurquer}.) To speed up? (\textbf{accélérer}.) To turn the wheel? (\textbf{braquer}.) To skid? (\textbf{déraper}.) To stall? (\textbf{caler}.) The pedal? (\textbf{la pédale}.) The windscreen? (\textbf{le pare-brise}.) The rear-view mirror? (\textbf{le rétroviseur}.) The horn? (\textbf{le klaxon}.) To honk? (\textbf{klaxonner}.) To mist up? (\textbf{embuer}.) The pump? (\textbf{la pompe}.) The diesel? (\textbf{le gazole}.) The fuel? (\textbf{le carburant}.) The tank? (\textbf{le réservoir}.) To start up? (\textbf{démarrer}.) The puncture? (\textbf{la crevaison}.) To get a flat tyre? (\textbf{crever}.) The tyre? (\textbf{le pneu}.) The jack? (\textbf{le cric}.) To pump up? (\textbf{gonfler}.) To get going again? (\textbf{dépanner}.) The tow truck? (\textbf{la dépanneuse}.) The breakdown mechanic? (\textbf{le dépanneur}.) To tow? (\textbf{remorquer}.) The garage owner? (\textbf{le garagiste}.) The bodywork? (\textbf{la carrosserie}.) The bonnet? (\textbf{le capot}.) The car door? (\textbf{la portière}.) The boot? (\textbf{le coffre}.) The back seat? (\textbf{la banquette}.) Parking? (\textbf{le stationnement}.) To park? (\textbf{stationner}.) The parking ticket machine? (\textbf{l'horodateur}.) The parallel parking space? (\textbf{le créneau}.) The car pound? (\textbf{la fourrière}.)
\end{tcolorbox}
